\subsection*{Kết luận chung}
\phantomsection\addcontentsline{toc}{subsection}{\numberline {}Kết luận chung}

Báo cáo đã xây dựng hoàn chỉnh bộ mô phỏng Monte Carlo cho HARQ trong 5G NTN và đánh giá hệ thống qua năm chiều hiệu năng. Bốn kết luận chính:

\begin{enumerate}
    \item \textbf{HARQ mang lại độ lợi link đáng kể ở LEO:} IR-HARQ với 4 TX đạt độ lợi 8{,}3 dB so với No HARQ tại BLER = $10^{-3}$ (MCS A, K = 15 dB), nhất quán với kết quả của~\cite{tuninato2025}. IR luôn vượt CC nhờ bất đẳng thức Jensen~(\ref{eq:jensen}).

    \item \textbf{Pipeline stall là rào cản kiến trúc không thể vượt qua với tiêu chuẩn hiện tại:} Chỉ 3/16 tổ hợp quỹ đạo--SCS có N$_{\min} \leq 32$. MEO và GEO bất khả thi ở mọi cấu hình SCS. Đây không phải hạn chế vật lý mà là hạn chế giao thức -- cần sửa đổi chuẩn 3GPP.

    \item \textbf{Tại GEO, vô hiệu hóa HARQ là quyết định đúng đắn:} RLC ARQ tăng SE Goodput 16{,}7 lần (0{,}030 $\to$ 0{,}500 bit/s/Hz) với độ trễ gần tương đương, xác nhận định lượng khuyến nghị của TR 38.821 Mục 6.4.2~\cite{tr38821}.

    \item \textbf{HARQ cải thiện mạnh hiệu quả năng lượng ở SNR thấp:} Năng lượng/bit với HARQ thấp hơn No HARQ $\sim 10^9$ lần tại $E_s/N_0 = -5$ dB -- ý nghĩa thực tiễn rõ ràng cho thiết bị IoT vệ tinh chạy bằng pin. Tuy nhiên, pipeline stall ở quỹ đạo cao làm tăng $E_{\text{bit}}$ tỷ lệ $1/\text{util}$ (GEO: 17 lần so với LEO\_{600}).
\end{enumerate}

\textbf{Khuyến nghị cấu hình:}

\begin{table}[H]
    \centering
    \caption[Cấu hình HARQ khuyến nghị theo quỹ đạo]{\bfseries\fontsize{12pt}{0pt}\selectfont
    Cấu hình HARQ tối ưu khuyến nghị cho từng quỹ đạo NTN}
    \label{tab:recommendation}
    \begin{tabularx}{0.98\textwidth}{
    | >{\centering\arraybackslash}X
    | >{\centering\arraybackslash}X
    | >{\centering\arraybackslash}X |}
    \hline
    \bfseries Quỹ đạo & \bfseries Cấu hình khuyến nghị & \bfseries Lý do \\\hline
    LEO 600 km & IR-HARQ, max\_retx=3, SCS $\leq$ 30 kHz & util=100\%, IR cho độ lợi tối đa \\\hline
    LEO 1200 km & IR-HARQ, max\_retx=3, SCS 15 kHz & Chỉ SCS 15 kHz khả thi \\\hline
    MEO 10000 km & HARQ disabled + RLC ARQ & N$_{\min}=95$--749, vượt giới hạn 32 \\\hline
    GEO 35786 km & \textbf{HARQ disabled + RLC ARQ} & SE tăng 16{,}7$\times$; N$_{\min}=272$--2166 \\\hline
    \end{tabularx}
\end{table}

\subsection*{Hướng phát triển}
\phantomsection\addcontentsline{toc}{subsection}{\numberline {}Hướng phát triển}

\begin{enumerate}
    \item \textbf{Kênh hai trạng thái LMS:} Mô hình Markov good/bad (TR 38.811 Mục 6.7.2~\cite{tr38811}) để xét hiệu ứng che khuất -- sẽ làm giảm hiệu quả HARQ trong trạng thái bad.
    \item \textbf{HARQ với N mở rộng:} Đánh giá lợi ích nếu 3GPP Release 19 tăng giới hạn N$_{\max}$ (Work Item NTN HARQ enhancement) hoặc cho phép HARQ bất đồng bộ.
    \item \textbf{Tương quan thời gian kênh:} Mô phỏng với kênh có tương quan thời gian (mô hình Jakes) để đánh giá lợi ích phân tập trong slow-fading (UE tĩnh, $T_c \gg T_{\text{slot}}$).
    \item \textbf{Kết hợp với AMC:} Phân tích link adaptation động -- chọn MCS theo trạng thái kênh dự đoán để giảm avg\_ntx và tối ưu E$_{\text{bit}}$.
\end{enumerate}

\subsection*{Kiến nghị và đề xuất}
\phantomsection\addcontentsline{toc}{subsection}{\numberline {}Kiến nghị và đề xuất}

Bộ mô phỏng này có thể được mở rộng trực tiếp để hỗ trợ các nghiên cứu tiếp theo:
\begin{itemize}
    \item Thêm module kênh hai trạng thái (src/channel/lms\_channel.py đã có cấu trúc sẵn).
    \item Tích hợp bộ mã hóa/giải mã LDPC thực (thư viện \texttt{sionna} hoặc \texttt{openairinterface}) để thay thế mô hình MI-threshold.
    \item Mở rộng Sim 04 với nhiều giá trị N (từ 32 đến 512) để tìm N ngưỡng cho MEO/GEO.
\end{itemize}
